\documentclass[12pt]{article}
\usepackage[a4paper,vmargin=2cm,hmargin=3cm]{geometry}
\usepackage[utf8]{inputenc}
\usepackage{setspace}
\usepackage[english]{babel}
\setlength{\parindent}{0em}
\usepackage{lmodern} 
\usepackage{amsmath}
\usepackage{amssymb}
\usepackage{graphicx}
\usepackage[export]{adjustbox}
\onehalfspacing
\usepackage{scalerel,stackengine}
\usepackage{float}
\usepackage[dvipsnames]{xcolor}
\usepackage{amsmath}
\DeclareMathOperator*{\argmax}{arg\,max}
\DeclareMathOperator*{\argmin}{arg\,min}
\newcommand{\indep}{\perp \!\!\! \perp}

\title{\textbf{Syllabus}}
\author{\textbf{Econometrics (25117)}}
\date{Universitat Pompeu Fabra, 2026-2027}

\begin{document}
\maketitle

\section*{Course Details}
\begin{itemize}
    \item \textbf{Lecturer}\\
    Juan Rubio-Ramirez
    \item \textbf{Academic Term}\\
    Fall Term (UPF regular course calendar): September 28 -- December 18, 2026
    \item \textbf{Sessions}\\
    18 Theory Sessions\\
    7 Seminar Sessions (starting in week 3)
    \item \textbf{Exams}\\
    Midterm 1 --- Date to be announced\\
    Midterm 2 --- Date to be announced\\
    Cumulative Final Exam --- Date to be announced
\end{itemize}
\section*{Course Description}
The course aims to introduce students to econometric methods and techniques. Econometrics lies at the intersection of economics and statistics, and econometric methods are widely used in applied work to estimate economic models, test hypotheses, and produce forecasts for future trends in relevant economic and financial variables.\\

The course is structured into two main blocks to provide a comprehensive understanding of econometric methods.

\begin{enumerate}
    \item Fundamental Inference Methods and OLS Regression --- In this initial block, students will learn essential statistical inference techniques that form the foundation of econometric analysis. They will be introduced to ordinary least squares (OLS) regression models, which are pivotal for estimating relationships between variables. This section will cover key concepts such as model specification, interpretation of coefficients, hypothesis testing, and the assumptions underlying OLS. Through worked examples and applications, students will gain hands-on experience in using OLS regression to analyze real-world data.
    \item Limitations of OLS and Alternative Models --- The second block addresses the limitations inherent in OLS regression models, particularly when dealing with more complex data structures or non-linear relationships. Students will explore various challenges such as multicollinearity, heteroscedasticity, and autocorrelation that can affect model accuracy. To enhance their analytical toolkit, the block will introduce alternative econometric models and data configurations.
\end{enumerate}

The teaching methodology is based on lectures, seminars, and individual study. The lectures aim to deepen students' understanding of the topics presented. They introduce various econometric estimators, explore their theoretical properties, and provide illustrative examples and applications to highlight the empirical relevance of the theory. During lectures, I will work through selected exercises in class. Seminars will provide additional practice and help students develop intuition for the key econometric concepts.\\

Additionally, students are expected to dedicate time to self-study by reading the book chapter references as well as the supplementary materials suggested in class and working through exercises to reinforce the skills learned during the lectures. This comprehensive approach ensures that students not only understand the theoretical aspects of basic econometrics but also become proficient in applying these methods in real-world scenarios.\\

By the end of the course, students should have an overall robust understanding of both foundational and basic econometric techniques, equipping them to tackle a wide range of economic and financial analysis challenges.


\section*{Grading}
To pass the course students have to obtain 50 out of 100 points according to the following distribution. Active participation in lectures and seminars is expected.
\begin{enumerate}
    \setlength{\itemsep}{0pt}
    \item Cumulative Final Exam: 60 points
    \item Midterm 2: 20 points
    \item Midterm 1: 15 points
    \item Participation: 5 points
\end{enumerate}

There will be a second chance to take the final exam at the end of the fall term only for those who failed to pass the course.


\section*{Requirements}
This course is intensive and requires a solid commitment to self-study. Students should be prepared to engage in readings and complete assignments and exercises. Active participation is encouraged, as it will enhance understanding and collaboration. A foundational background in probability, statistics, and economics is required for engaging with the course material. Students are expected to read the primary textbook references and any supplementary materials if need be.

\section*{Contents}

The core structure of the course follows \textit{Introduction to Econometrics, 4th Edition, Global Edition} by Stock and Watson. The selected chapters from this book are \underline{highly} recommended readings. Students may use the print book, library access, or the authorized Pearson eText version; no official downloadable PDF version is required for the course.

\begin{itemize}
    \item \textbf{Lecture I: Introduction and Review of Probability}
    \begin{itemize}
        \item Introduction to Econometrics, 4th Edition, Global Edition, by Stock and Watson -- Chapter 1-2

    \end{itemize}

    \item \textbf{Lecture II: Review of Statistics}
    \begin{itemize}
        \item Introduction to Econometrics, 4th Edition, Global Edition, by Stock and Watson -- Chapter 3\\[1em]
    \end{itemize}
    
    \item \textbf{Lecture III: Simple Linear Regression}
    \begin{itemize}
        \item Introduction to Econometrics, 4th Edition, Global Edition, by Stock and Watson -- Chapter 4
    \end{itemize}
    
    \item \textbf{Lecture IV: OLS Implementation}
    \begin{itemize}
        \item Introduction to Econometrics, 4th Edition, Global Edition, by Stock and Watson -- Chapter 5
    \end{itemize}
    
     \item \textbf{Lecture V: Multivariate Linear Regression}
     \begin{itemize}
        \item Introduction to Econometrics, 4th Edition, Global Edition, by Stock and Watson -- Chapter 6
    \end{itemize}
    
     \item \textbf{Lecture VI: Inference and Interpretation}
     \begin{itemize}
        \item Introduction to Econometrics, 4th Edition, Global Edition, by Stock and Watson -- Chapter 7
    \end{itemize}

    \item \textbf{Lecture VII: Functional Forms}
    \begin{itemize}
        \item Introduction to Econometrics, 4th Edition, Global Edition, by Stock and Watson -- Chapter 8
    \end{itemize}

    \item \textbf{Lecture VIII: Threats to Identification}
    \begin{itemize}
        \item Introduction to Econometrics, 4th Edition, Global Edition, by Stock and Watson -- Chapter 9

    \end{itemize}

    \item \textbf{Lecture IX: Instrumental Variables}
    \begin{itemize}
        \item Introduction to Econometrics, 4th Edition, Global Edition, by Stock and Watson -- Chapter 12

    \end{itemize}

    \item \textbf{Lecture X: Panel Data}
    \begin{itemize}
        \item Introduction to Econometrics, 4th Edition, Global Edition, by Stock and Watson -- Chapter 10-13
    \end{itemize}

    \item \textbf{Lecture XI: Binary Outcomes}
    \begin{itemize}
        \item Introduction to Econometrics, 4th Edition, Global Edition, by Stock and Watson -- Chapter 11

    \end{itemize}

\end{itemize}




\end{document}
